\section{Proofs for integer recourse}\label{app:int}

This appendix proves the results of \cref{sec:int}. \Cref{app:int:solver}
proves the exact box solver, \cref{app:int:native} the native integer
theorems and the oracle statements, \cref{app:int:flows} the flow results,
\cref{app:int:tu} the totally unimodular extension, and
\cref{app:int:strong-field,app:int:lowrank} the strong-noise and low-rank
results. The tools \ref{it:app:rec:local}--\ref{it:app:rec:faces} of
\cref{app:rec:tools} are used freely. We also use the following univariate
facts \citep{basu2006-algorithms-in-real-algebraic-geometry,mehlhorn2015-from-approximate-factorization-to-root}: for integer
polynomials of degree at most $\delta$ and coefficient bit length at most
$\tau$, greatest common divisors, exact quotients, squarefree parts, modular
inverses, resultants, isolation of all real roots by disjoint rational
intervals, the sign of another such polynomial at an isolated root, and
refinement of an isolating interval to width $2^{-q}$ take
$\poly(\delta,\tau,q)$ bit operations with an absolute exponent. Algebraic
degrees are polynomial in $\delta$, and intermediate bit lengths are
polynomial in $\delta,\tau,q$. Every nonzero root $t$ of an integer polynomial
$\sum_ja_jt^j$ with $a_0\ne0$ satisfies
$\abs t\ge\abs{a_0}/(\abs{a_0}+\max_{j\ge1}\abs{a_j})$, by Cauchy's bound applied
to the reversed polynomial.

\subsection{The exact box solver}\label{app:int:solver}

Let $f$ and $B=\prod_{i=1}^k[\ell_i,u_i]$ be as in \cref{thm:int:solver}.
Substitute the coordinates with $\ell_i=u_i$; then every interval is
nondegenerate. A \emph{face} $\Phi$ is a partition of $[k]$ into sets
$\mathsf L,\mathsf U,\mathsf F$ of coordinates fixed at the lower bound,
fixed at the upper bound, and free. Write $r=\abs{\mathsf F}$ and
$f_\Phi(y)$, $y\in\R^r$, for $f$ with the fixed coordinates substituted; its
degree is at most $d$. Let $D$ be the least even integer greater than $d$,
let $a=D-1\ge d$, and let $z$ be an indeterminate. For a face with $r\ge1$
put
\begin{equation}\label{eq:app:int:system}
 g_i(y)=y_i^a+\frac zD\,\partial_if_\Phi(y)\qquad(i\in\mathsf F).
\end{equation}
For a real $z=1/\varepsilon>0$ these are the stationary equations of
$f_\Phi(y)+\varepsilon\sum_{i\in\mathsf F}y_i^D$ divided by $\varepsilon D$, and
$f+\varepsilon\sum_{i=1}^kx_i^D$ restricted to the face is this function plus a
constant. Let $\mathcal M=\{y^e:0\le e_i<a\}$ and $N=\abs{\mathcal M}=a^r$.

\paragraph{The algorithm.}
For every face: if $r=0$, the vertex is a candidate. If $r\ge1$:
\begin{enumerate}[label=(\arabic*)]
\item Compute the normal forms of all monomials of total degree at most
$R=r(a-1)+1$ with respect to \eqref{eq:app:int:system}, in order of increasing
degree: a monomial in $\mathcal M$ is its own normal form, and otherwise
$\mathrm{NF}(y_i^ay^{e'})=-\frac zD\sum_\varrho c_\varrho\mathrm{NF}(y^\varrho y^{e'})$,
where $\partial_if_\Phi=\sum_\varrho c_\varrho y^\varrho$. Read off the matrices
$\mathsf X_i(z)\in\Q[z]^{N\times N}$ of multiplication by $y_i$ in the basis
$\mathcal M$.
\item For $t=0,1,\ldots,Q_r-1$, $Q_r=1+(r-1)(N+\binom N2)$, let
$\lambda=(1,t,\ldots,t^{r-1})$ and $\mathsf X_\lambda=\sum_i\lambda_i\mathsf X_i$. Compute
$\chi(z,T)=\det(TI-\mathsf X_\lambda)$ and, for $i\in\mathsf F$,
$\chi_i(z,T)=-\partial_u\det(TI-\mathsf X_\lambda-u\mathsf X_i)\vert_{u=0}$, by interpolation in
$(z,T,u)$.
\item Let $K_\lambda=\deg_z\chi$ and $p=[z^{K_\lambda}]\chi\in\Q[T]$. Skip the
form if $p$ is constant, if some $\deg_z\chi_i>K_\lambda$, if
$G=\gcd(p,p')$ does not divide $b_i=[z^{K_\lambda}]\chi_i$ for some $i$, or if
$\gcd(p'/G,p/G)\ne1$. Otherwise let $P=p/G$, $A=p'/G$, and $B_i=b_i/G$.
Multiply $P$ by a positive common denominator of its coefficients so that
$P\in\Z[T]$, and put $r_i=B_iA^{-1}\bmod P$. This scaling changes neither
the roots nor the quotient ideal.
\item For every real root $\theta$ of $P$, keep the point with free coordinates
$r_i(\theta)$ and the fixed coordinates of the face if it lies in $B$; the
membership test consists of signs of $r_i-\ell_i$ and $u_i-r_i$ at $\theta$.
\end{enumerate}
Among all kept points and vertices, select the first one of least value in
a fixed order, compared exactly by \cref{lem:int:values}. Return its integer
polynomial $P$, the isolating interval
of its root $\theta$, the maps $r_i$ of the free coordinates and constant maps
for the fixed ones, the value map $r_f$ of \cref{lem:int:values}, and, as the
auxiliary representation, the value polynomial of that lemma with an isolating
interval of the value. A selected vertex is returned with $P(T)=T$,
$\theta=0$, and constant maps.

\begin{lemma}[Quotient algebra]\label{lem:int:quotient}
Let $\mathbb F$ be a field containing $\Q$ and $\zeta\in\mathbb F$, or let
$\mathbb F$ contain $\Q(z)$ and $\zeta=z$. Under every degree-compatible monomial
order, the polynomials $g_i$ with $z=\zeta$ form a Gr\"obner basis of the ideal
$\mathcal I_\zeta\subseteq\mathbb F[y]$ they generate, $\mathcal M$ is a basis
of $\mathbb F[y]/\mathcal I_\zeta$, and the matrices of multiplication by $y_i$
in this basis are $\mathsf X_i(\zeta)$. They commute.
\end{lemma}

\begin{proof}
The term $y_i^a$ has coefficient one and degree $a$, and every other term of
$g_i$ has degree at most $d-1<a$. So the leading monomials are $y_i^a$
whatever $\zeta$ is, and they are pairwise coprime. By Buchberger's first
criterion all S-polynomials reduce to zero, so the $g_i$ form a Gr\"obner basis,
and the monomials not divisible by any $y_i^a$, namely $\mathcal M$, form a
basis of the quotient \citep{cox2015-ideals-varieties-and-algorithms}. Reduction by the $g_i$
divides only by the leading coefficient one, so the normal form of a monomial
is the specialization at $\zeta$ of the normal form computed over $\Q[z]$.
Each recursion step of the algorithm lowers the total degree by at least
$a-(d-1)\ge1$, so the normal forms of all monomials of degree at most $R$,
which include $y_i\mathfrak m$ for $\mathfrak m\in\mathcal M$, are computed. Multiplication
operators of a commutative algebra commute.
\end{proof}

Let $\mathcal P$ be the field of Puiseux series
$\bigcup_{m\ge1}\overline\Q((\varepsilon^{1/m}))$; it is algebraically closed
\citep{basu2006-algorithms-in-real-algebraic-geometry}. Embed $\Q(z)$ in $\mathcal P$ by $z\mapsto\varepsilon^{-1}$.
For $s\in\mathcal P$ let $\operatorname{val}(s)$ be its least exponent, and for a
vector or a polynomial in $T$ the least valuation of its entries or
coefficients.

\begin{lemma}[Joint eigenvalues]\label{lem:int:triangular}
There are points $\xi^{(1)},\ldots,\xi^{(N)}\in\mathcal P^r$, not necessarily
distinct, such that every $\xi^{(t)}$ is a common zero of the $g_i$ (with
$z=\varepsilon^{-1}$) and $\det\bigl(TI-\sum_i\lambda_i\mathsf X_i(\varepsilon^{-1})\bigr)
=\prod_{t=1}^N(T-\lambda^{\mathsf T}\xi^{(t)})$ for every $\lambda\in\mathcal P^r$.
\end{lemma}

\begin{proof}
Commuting matrices over an algebraically closed field are simultaneously
upper triangularizable \citep{horn2013-matrix-analysis}. Let $\Pi$ triangularize all
$\mathsf X_i(\varepsilon^{-1})$, and let $\xi^{(t)}_i$ be the $t$th diagonal entry of
$\Pi^{-1}\mathsf X_i\Pi$. Then $\Pi^{-1}\mathsf X_\lambda\Pi$ is triangular with diagonal
entries $\lambda^{\mathsf T}\xi^{(t)}$, which gives the determinant. For a
polynomial $h$, let $\mathsf X_h$ be the matrix of multiplication by $h$; the map
sending $h$ to the $t$th diagonal entry of $\Pi^{-1}\mathsf X_h\Pi$ is a ring
homomorphism, because diagonal entries of sums and products of upper triangular
matrices are sums and products of diagonal entries. It sends $y_i$ to
$\xi^{(t)}_i$ and every element of the ideal to $0$. So $g_i(\xi^{(t)})=0$.
\end{proof}

Call $t$ \emph{bounded} if $\operatorname{val}(\xi^{(t)})\ge0$, and let
$v^{(t)}\in\overline\Q^r$ be the vector of its $\varepsilon^0$-coefficients, its
\emph{limit}. Otherwise let $s_t=-\operatorname{val}(\xi^{(t)})>0$ and let
$w^{(t)}\ne0$ be the vector of $\varepsilon^{-s_t}$-coefficients, its
\emph{leading pole}. Let $\mathcal L$ be the set of distinct limits, $\mu_v\ge1$
the number of bounded $t$ with limit $v$, $\kappa=\sum_{t\text{ unbounded}}s_t$,
and $C(\lambda)=\prod_{t\text{ unbounded}}(-\lambda^{\mathsf T}w^{(t)})$, a nonzero
polynomial. Write
$\chi(\lambda,z,T)=\det(TI-\sum_i\lambda_i\mathsf X_i(z))=\sum_mc_m(\lambda,T)z^m$ with
$c_m\in\Q[\lambda,T]$.

\begin{lemma}[The leading coefficient]\label{lem:int:leading}
The number $\kappa$ is an integer, $c_m=0$ for $m>\kappa$, and
$c_\kappa(\lambda,T)=C(\lambda)\prod_{v\in\mathcal L}(T-\lambda^{\mathsf T}v)^{\mu_v}$
as polynomials in $(\lambda,T)$.
\end{lemma}

\begin{proof}
Let $\lambda\in\overline\Q^r$ satisfy $\lambda^{\mathsf T}w^{(t)}\ne0$ for every
unbounded $t$. For unbounded $t$, $\lambda^{\mathsf T}\xi^{(t)}$ has valuation
$-s_t$ and leading coefficient $\lambda^{\mathsf T}w^{(t)}$, so
$T-\lambda^{\mathsf T}\xi^{(t)}=(-\lambda^{\mathsf T}\xi^{(t)})
(1-T/\lambda^{\mathsf T}\xi^{(t)})$, where the second factor is $1$ plus a
polynomial of positive valuation. For bounded $t$,
$T-\lambda^{\mathsf T}\xi^{(t)}$ is $T-\lambda^{\mathsf T}v^{(t)}$ plus a
polynomial of positive valuation. By \cref{lem:int:triangular},
\[
 \varepsilon^\kappa\chi(\lambda,\varepsilon^{-1},T)
 =C(\lambda)\prod_{v\in\mathcal L}(T-\lambda^{\mathsf T}v)^{\mu_v}+\Omega(T),
 \qquad\operatorname{val}(\Omega)>0 .
\]
The left side equals $\sum_mc_m(\lambda,T)\varepsilon^{\kappa-m}$. Comparing the
coefficients of each power of $T$ shows that $c_m(\lambda,\cdot)=0$ when
$\kappa-m<0$, and that the $\varepsilon^0$-part of the left side equals the
nonzero first term on the right; hence some $m$ equals $\kappa$, and
$c_\kappa(\lambda,T)=C(\lambda)\prod_v(T-\lambda^{\mathsf T}v)^{\mu_v}$. These
identities hold on the complement of finitely many hyperplanes in
$\overline\Q^r$, so they hold as polynomial identities.
\end{proof}

Call a form $\lambda\in\Q^r$ \emph{good} for the face if
$\lambda^{\mathsf T}w^{(t)}\ne0$ for every unbounded $t$ and
$\lambda^{\mathsf T}v\ne\lambda^{\mathsf T}v'$ for distinct $v,v'\in\mathcal L$.

\begin{lemma}[Recovery from one root]\label{lem:int:recovery}
If $\lambda$ is good and $\mathcal L\ne\emptyset$, the form passes step~(3), and
for every $v\in\mathcal L$ the number $\theta=\lambda^{\mathsf T}v$ is a root of
$P$ with $r_i(\theta)=v_i$ for all $i$. If $\theta$ is real, $v$ is real.
\end{lemma}

\begin{proof}
By \cref{lem:int:leading}, $K_\lambda=\kappa$ because $C(\lambda)\ne0$, and
$p=C(\lambda)\prod_v(T-\alpha_v)^{\mu_v}$ with distinct $\alpha_v=\lambda^{\mathsf T}v$.
Since $\chi(\lambda+ue_i,z,T)=\det(TI-\mathsf X_\lambda-u\mathsf X_i)$, we have
$\chi_i=-\partial_{\lambda_i}\chi$; because $c_m\equiv0$ for $m>\kappa$, also
$\deg_z\chi_i\le\kappa$ and $b_i=-\partial_{\lambda_i}c_\kappa$. Differentiating
the identity of \cref{lem:int:leading},
\[
 b_i=-\frac{\partial C}{\partial\lambda_i}\prod_u(T-\alpha_u)^{\mu_u}
 +C(\lambda)\sum_v\mu_vv_i\,(T-\alpha_v)^{\mu_v-1}\prod_{u\ne v}(T-\alpha_u)^{\mu_u},
\]
while $p'=C(\lambda)\sum_v\mu_v(T-\alpha_v)^{\mu_v-1}\prod_{u\ne v}(T-\alpha_u)^{\mu_u}$
and $G=\gcd(p,p')$ is a constant times $\prod_v(T-\alpha_v)^{\mu_v-1}$. So $G$
divides $b_i$, $P$ is a constant times $\prod_v(T-\alpha_v)$, and with a common
constant $\kappa'\ne0$,
\[
 A(\alpha_v)=\kappa'\mu_v\prod_{u\ne v}(\alpha_v-\alpha_u)\ne0,\qquad
 B_i(\alpha_v)=\kappa'\mu_vv_i\prod_{u\ne v}(\alpha_v-\alpha_u),
\]
because the first term of $b_i$ keeps a factor $T-\alpha_v$ after division by
$G$, and $\mu_v\ne0$ in characteristic zero. Thus $\gcd(A,P)=1$ and
$r_i(\alpha_v)=B_i(\alpha_v)/A(\alpha_v)=v_i$. The polynomials $A$ and $B_i$ have
rational coefficients, so a real $\theta$ gives a real $v$.
\end{proof}

\begin{lemma}[Limits of real critical points]\label{lem:int:limits}
Let $\varepsilon_m\downarrow0$, and let $y^{(m)}\in\R^r$ be common zeros of the
$g_i$ with $z=1/\varepsilon_m$ that converge to $y^*$. Then $y^*\in\mathcal L$.
\end{lemma}

\begin{proof}
Let $\lambda\in\Q^r$. The row vector $e=(\mathfrak m(y^{(m)}))_{\mathfrak m\in\mathcal M}$ is
nonzero, because $1\in\mathcal M$, and satisfies
$e\mathsf X_i(1/\varepsilon_m)=y^{(m)}_ie$: evaluation at a common zero is well defined
on the quotient. So $\lambda^{\mathsf T}y^{(m)}$ is an eigenvalue of
$\mathsf X_\lambda(1/\varepsilon_m)$, and by \cref{lem:int:leading}
\[
 0=\varepsilon_m^\kappa\chi(\lambda,1/\varepsilon_m,\lambda^{\mathsf T}y^{(m)})
 =\sum_{m'\le\kappa}c_{m'}(\lambda,\lambda^{\mathsf T}y^{(m)})\varepsilon_m^{\kappa-m'} .
\]
The arguments are bounded, so letting $m\to\infty$ gives
$c_\kappa(\lambda,\lambda^{\mathsf T}y^*)=0$, that is,
$C(\lambda)\prod_{v\in\mathcal L}(\lambda^{\mathsf T}(y^*-v))^{\mu_v}=0$ for all
$\lambda\in\Q^r$. The polynomial on the left vanishes on the Zariski-dense set
$\Q^r$, so it is zero; $C\ne0$ and the polynomial ring is a domain, so some
linear factor $\lambda^{\mathsf T}(y^*-v)$ is identically zero. Hence $y^*=v$.
If $\mathcal L$ were empty the same argument would give $C\equiv0$.
\end{proof}

\begin{lemma}[Good forms]\label{lem:int:forms}
Among $\lambda(t)=(1,t,\ldots,t^{r-1})$, $t=0,\ldots,Q_r-1$, at least one form is
good.
\end{lemma}

\begin{proof}
A form fails for $t$ if $\lambda(t)^{\mathsf T}w=0$ for one of at most $N$
leading poles $w$, or $\lambda(t)^{\mathsf T}(v-v')=0$ for one of at most
$\binom N2$ pairs of distinct limits. Each condition says that a nonzero
polynomial in $t$ of degree at most $r-1$, with coefficient vector $w$ or
$v-v'$, vanishes; it has at most $r-1$ roots. So at most
$(r-1)(N+\binom N2)<Q_r$ integers fail.
\end{proof}

\begin{proof}[Proof of correctness in \cref{thm:int:solver}]
For $\varepsilon>0$ let $x(\varepsilon)$ minimize $f+\varepsilon\sum_ix_i^D$ on
$B$. Along a sequence $\varepsilon_m\downarrow0$ choose a convergent subsequence
with limit $\hat x$; since the perturbation converges to zero uniformly on $B$,
$\hat x$ minimizes $f$. Pass to a further subsequence on which every
$x(\varepsilon_m)$ lies in the relative interior of one face $\Phi$. If $\Phi$ is
a vertex, $\hat x$ is that vertex, a candidate. Otherwise the free part of
$x(\varepsilon_m)$ is an interior minimizer of the deformed face function, hence
a real common zero of \eqref{eq:app:int:system} with $z=1/\varepsilon_m$. By
\cref{lem:int:limits} its limit lies in $\mathcal L_\Phi$; by
\cref{lem:int:forms,lem:int:recovery} some enumerated form recovers it exactly
from a real root of $P$, and it lies in $B$. So $\hat x$ is kept. A bad form
that passes step~(3) still yields well-defined maps, because $A$ is invertible
modulo $P$, but its reconstructed points need not be critical limits; for
instance, its leading coefficient can retain the bounded projection of a
branch that escapes to infinity. Step~(4) keeps such a point only if it lies
in $B$. Every kept point is therefore feasible and has value at least
$\min_Bf$, so a kept point of least value is a global minimizer.
\end{proof}

\begin{lemma}[Values and comparisons]\label{lem:int:values}
For a kept point with root $\theta$ of $P$ and coordinates $r_i(\theta)$ (fixed
coordinates substituted), let $r_f=f(r(T))\bmod P$ be the remainder of
$f(r_1(T),\ldots,r_k(T))$ on division by $P$. Then $\deg r_f<\deg P$ and
$r_f(\theta)=f(r(\theta))$, the value of the point. Write $r_f=h/q_0$ with
$h\in\Z[T]$ and $q_0\in\Z\setminus\{0\}$. The resultant
$\operatorname{Res}_T(P(T),q_0W-h(T))\in\Z[W]$ is nonzero, has degree at most
$\deg P$, and has the value as a root, which is identified by refining the
interval of $\theta$. Two such values are compared exactly, including
equality, by isolating the real roots of the squarefree part of the product
of their two value polynomials and locating both values among them.
\end{lemma}

\begin{proof}
$P$ divides $f(r(T))-r_f(T)$, and $P(\theta)=0$, so $r_f(\theta)=f(r(\theta))$.
Up to a nonzero constant factor, the resultant is
$\prod_\vartheta(q_0W-h(\vartheta))$ over the roots $\vartheta$ of $P$, a
nonzero polynomial of degree $\deg P$ with the root $h(\theta)/q_0$. Refining
the interval of $\theta$ encloses $h(\theta)/q_0$ by interval evaluation; by the
root separation bound of the squarefree part, polynomially many bits suffice
to single out its root. Equal values are the same root of the squarefree
product, so ties are decided exactly.
\end{proof}

The value map $r_f$ belongs to the output, so the optimizer and the value are
polynomials in the same root, as \cref{def:model:outputs}\ref{it:model:algebraic}
requires. The value polynomial and its isolating interval are an auxiliary
representation for comparisons; composing them with the coordinate
representation is never needed.

\begin{lemma}[Cost]\label{lem:int:cost}
The algorithm uses $c_d^k\poly_d(H)$ bit operations for a constant $c_d$ that
depends only on $d$. Its output, and the representations of the value, have
length $c_d^k\poly_d(H)$, and $q$-bit refinements cost $c_d^k\poly_d(H+q)$.
\end{lemma}

\begin{proof}
Fix a face with $r\ge1$ free coordinates. There are
$\binom{r+R}r\le2^{r+R}\le2^{Dr}$ monomials of degree at most $R$ and
$N=a^r\le2^{Dr}$. Clear the denominators of $\partial_if_\Phi/D$ by one integer of
$O_d(H)$ bits. Each recursion step multiplies by $z$ and by integer coefficients
of $O_d(H)$ bits; a normal form is a sum over reduction paths of at most $R$
steps, of which there are at most $(s+1)^R$, where $s=\poly_d(r)$ bounds the
number of monomials of a derivative. So every normal form has $z$-degree at most
$R$ and coefficients of $O_d(R(H+\log(s+1)))$ bits; memoization computes all of
them without enumerating paths, in $2^{O_d(r)}\poly_d(H)$ operations. The forms
have $O(r^2\log(N+1))$-bit entries. Each determinant has order $N$, degree at
most $N$ in $T$ and $u$ and at most $RN$ in $z$, and is interpolated from
$(RN+1)(N+1)^2$ rational determinants of polynomial length; so $\chi$ and the
$\chi_i$ have $\poly(N,H,r)$-bit coefficients. The $Q_r=O(rN^2)$ forms, the
univariate steps (3)--(4) on polynomials of degree at most $N$, the at most $N$
real roots per form, and the value and comparison steps of
\cref{lem:int:values} each cost $\poly(N,H,r)$; the value map is obtained from
at most $H$ monomials of degree at most $d$, each a product of at most $d$
maps reduced modulo $P$. Summing over at most $3^k$ faces gives
$c_d^k\poly_d(H)$ with $c_d$ depending on $D$ only.

\emph{Refinement.} Refine the interval of $\theta$ and evaluate the $r_i$ and
$r_f$ by interval arithmetic with derivative bounds on the interval;
polynomially many additional bits give every coordinate within
$2^{-q-\lceil\log_2(k+1)\rceil}$, hence a point within Euclidean distance
$2^{-q}$, and the value within $2^{-q}$. Clipping a rational approximation to
$B$ does not increase its distance to the exact point, which lies in $B$. With
a rational bound $G_f\ge\max\{1,\sup_B\norm{\nabla f}_1\}$ from monomial sums,
coordinates within $2^{-q}/(2G_f)$ give a feasible rational point whose value
exceeds the optimum by at most $2^{-q-1}$; enclosing the value with width
$2^{-q-1}$ gives a lower bound, and so a certified gap $2^{-q}$. The value
gap alone would not bound the distance to $x^*$; the two refinements are
separate.
\end{proof}

\begin{lemma}[An explicit constant]\label{lem:int:budget}
Implement the algorithm with schoolbook integer arithmetic, fraction-free
Bareiss determinants, tensor interpolation on consecutive integer nodes,
subresultant greatest common divisors and B\'ezout coefficients, and squarefree
Sturm sequences for isolation, sign determination, and refinement. Then its bit
work for dimension $k$, input length $H$, and $q$ refinement bits is at most
$C_d\,2^{10000Dk}(H+q+1)^{10000}$ for a constant $C_d$ depending only on $d$. In
particular $c_d=2^{10000D}$ is a valid constant in \cref{thm:int:solver} and,
with a larger $C_d$, in \cref{cor:int:mixed-solver}.
\end{lemma}

\begin{proof}
Put $U=2^{Dk}(H+q+d+2)$. The ledger is deliberately loose. On a face with
$r\le k$ free coordinates the number of memoized monomials and $N$ are at most
$2^{Dk}$; normal forms have $z$-degree $O_d(k)$ and coefficient bit length
$O_d((H+k+1)^2)$; there are $O(rN^2)$ forms with $O_d(k^2)$-bit coefficients;
determinants have order $N$ and degrees $N$, $N$, $O_d(kN)$, so their
interpolation grids have $O_d(U^4)$ nodes, sampled values have bit length
$O_d(U^5)$, and interpolated coefficients $O_d(U^8)$. Subresultant and Sylvester
bounds keep greatest common divisors, inverses, and the $r_i$ within degree and
bit length $O_d(U^{12})$; composition with $f$ and the value resultant stay
within $O_d(U^{18})$; Cauchy and squarefree separation bounds, derivative bounds
for evaluation, and the $q$ extra bits require precision $O_d(U^{25})$. Every
primitive call therefore has input length at most $C_dU^{50}$, and the number of
primitive calls, over at most $3^k\le2^{Dk}$ faces, forms, coordinates, nodes,
roots, and comparisons, is at most $C_dU^{10}$. Each fixed elementary routine
costs at most $A(\ell+2)^{100}$ bit operations on inputs of length $\ell$, for an
absolute constant $A$: schoolbook multiplication and division are quadratic;
Bareiss elimination of order $N\le\ell$ performs $O(N^3)$ exact operations on
integers of $O(\ell^2)$ bits; Lagrange interpolation on $\nu\le\ell$ integer
nodes performs $O(\nu^2)$ operations on numbers of $O(\ell^2)$ bits; and
subresultant sequences, B\'ezout coefficients, and Sturm sequences of
polynomials of degree $\delta\le\ell$ perform $O(\delta^2)$ operations on
integers of $O(\ell^2)$ bits \citep{basu2006-algorithms-in-real-algebraic-geometry}, while isolation, sign
determination at an isolated root, and refinement by bisection evaluate
such sequences at $O(\ell^2)$ rational points of $O(\ell^2)$ bits. Each of
these bounds is at most $A(\ell+2)^{100}$. The total is at most
$C'_dU^{5010}\le C_d2^{10000Dk}(H+q+1)^{10000}$.
\end{proof}

The constant is an explicit proved bound used to define noise laws and fallback
budgets before sampling. A smaller constant with a proved bound of the same form
may replace it.

\begin{proof}[Proof of \cref{cor:int:mixed-solver}]
Enumerate the integer labels, substitute each, and apply \cref{thm:int:solver};
substitution keeps the degree and gives coefficients of length $\poly_d(H)$.
Keep the current best value and compare each new value with it by
\cref{lem:int:values}. Each comparison involves two univariate polynomials of
degree at most $(D-1)^k$, so the total is $c_d^k\prod_iN_i\poly_d(H)$ after
enlarging $c_d$ by a factor depending only on $d$. Only the winner's
representation is retained; isolating its roots afresh in their own polynomials
removes any extra precision used in comparisons.
\end{proof}

\subsection{Native integer recourse and integer oracles}\label{app:int:native}

Throughout, $G\ge\max\{1,\sup\sum_{i\in C}\abs{\partial_{v_i}F_\gamma}\}$ over
$[0,1]^k\times\prod_i[\ell_i,u_i]$ and all noise in $[-\sigma,\sigma]^n$, computed
by monomial bounds; it satisfies the hypothesis of \cref{lem:rec:value}(c).

\begin{lemma}[Competing labels]\label{lem:int:label}
Let $\mathcal H$ be a core box, $c\in\mathcal H$, and $z\in Y$ with
$F_\gamma(c,z)=V(c)$. The sets $Y_i^-=\{z'\in Y:z'_i\le z_i-1\}$ and
$Y_i^+=\{z'\in Y:z'_i\ge z_i+1\}$ cover $Y\setminus\{z\}$. If
\eqref{eq:int:label} holds with $V_{\rm other}(c)$ the least exact value over
the nonempty ones ($+\infty$ if none), then $z$ is the unique minimizer of
$F_\gamma(v,\cdot)$ on $Y$ for every $v\in\mathcal H$. If moreover $\mathcal H$
contains the core of every global minimizer, every global minimizer lies in
$\mathcal H\times\{z\}$ and $\min_{v\in[0,1]^k}F_\gamma(v,z)=f^*$.
\end{lemma}

\begin{proof}
If $Y=\{z\}$, the conclusions are immediate. Otherwise,
a label $z'\ne z$ differs in some coordinate, and by integrality $z'_i\le z_i-1$
or $z'_i\ge z_i+1$. For $v\in\mathcal H$, \cref{lem:rec:value}(c), applied to
$Y\setminus\{z\}$ and to $\{z\}$, gives
\begin{align*}
 V_{Y\setminus\{z\}}(v)&\ge V_{\rm other}(c)-Gw_\infty(\mathcal H)\\
 &>V(c)+Gw_\infty(\mathcal H)\\
 &=F_\gamma(c,z)+Gw_\infty(\mathcal H)\ge F_\gamma(v,z).
\end{align*}
A global minimizer $(a,z^*)$
with $a\in\mathcal H$ therefore has $z^*=z$. The slice
$[0,1]^k\times\{z\}$ lies in $X$ and contains $(a,z)$, so its minimum is $f^*$.
\end{proof}

\begin{lemma}[Stopping for labels]\label{lem:int:native-stop}
Suppose $F_\gamma$ has a unique global minimizer $x^*=(a,a_R)$ with point growth
at least $g_0$ on $X$, let $A_0=2+kL/g_0$, and let $h=2^{-j}$ satisfy
\[
 h\le\min\Bigl\{\frac1{4A_0},\ \frac{g_0}{16GA_0}\Bigr\}.
\]
Then the label test at level $j$ passes, with $z_j=a_R$.
\end{lemma}

\begin{proof}
By \cref{prop:rec:search} in exact mode, every retained cell has a corner $v$
with $V(v)\le f^*+2E_j$, and growth gives $\norm{v-a}\le h\sqrt{kL/(4g_0)}$; so
$\mathcal H_j\subseteq a+[-A_0h,A_0h]^k$ and $w_\infty(\mathcal H_j)\le2A_0h$. The
incumbent has $F_\gamma(c_j,z_j)\le f^*+E_j$, so
$\norm{(c_j,z_j)-x^*}\le\sqrt{E_j/g_0}\le A_0h\le\frac14$, which forces
$z_j=a_R$ because distinct labels are at distance at least one. Every other
label $z'$ has $F_\gamma(c_j,z')-f^*\ge g_0\norm{z'-a_R}^2\ge g_0$, so
$V_{\rm other}(c_j)-V(c_j)\ge g_0-E_j\ge\frac78g_0$, since
$E_j=kLh^2/8\le g_0/256$. Finally $2Gw_\infty(\mathcal H_j)\le4GA_0h\le g_0/4$.
\end{proof}

\begin{proof}[Proof of \cref{thm:int:native}]
\emph{Base choices.} Let $B=\max\{2,N_ZC_dc_d^k\}$, where $C_d,c_d$ are those of
\cref{lem:int:budget} enlarged by \cref{cor:int:mixed-solver}, so that enumerating
all labels, checking feasibility, solving each core slice, and comparing values
costs at most $B(I+\log M+1)^{e}$ for a fixed exponent $e$. Let
$S=k+\sum_i(u_i-\ell_i)$, $n=k+n_R$, and let $C_{\rm tail}$ be the constant of
\ref{it:app:rec:finite} for $\mathcal D=X$, described by the box, the rows of
$Az\le b$, and one disjunction per integer coordinate, with all $n$ coordinates
perturbed; $\log C_{\rm tail}\le\poly_d(I)$. Choose $\rho=1/(4B)$,
$g_0=\rho\sigma/(2S)$, $A_0=2+kL/g_0$, the least $J\ge0$ with
$2^{-J}\le\min\{1/(4A_0),g_0/(16GA_0)\}$, and the least power of two
$M\ge\max\{2,2^J,4nC_{\rm tail}/\rho\}$. All have polynomial bit length, and
$J,\log_2M\le\poly_d(I)$.

\emph{Algorithm and correctness.} At levels $j=0,\ldots,J$ run the core search in
exact mode with the integer oracle, and test \eqref{eq:int:label} at
$(c_j,z_j)$ with at most $2n_R$ further oracle calls on tightened bounds. When it
passes, apply \cref{thm:int:solver} to $F_\gamma(\cdot,z_j)$ on $[0,1]^k$ and
return. Otherwise, after level $J$, enumerate the labels on the same draw as in
\cref{cor:int:mixed-solver}. By \cref{prop:rec:search}(a),(b) and
\cref{lem:int:label}, every return is correct.

\emph{Probability.} By \cref{lem:int:native-stop}, the fallback runs only if
$g_*<g_0$, which by \ref{it:app:rec:finite} has probability at most
$Sg_0/\sigma+2nC_{\rm tail}/M\le\rho=1/(4B)$.

\emph{Work and output.} By \cref{prop:rec:search} the expected number of core
calls is at most $2^k+8^kJQ_{\rm ex}$; each call, including the label calls,
has polynomial query length and costs $\poly_d(I)$. The core slice has
coefficients of length $\poly_d(I+\log M)=\poly_d(I)$, so the completion costs
$c_d^k\poly_d(I)$ and returns a representation of that length. The fallback
contributes $\poly_d(I)$ in expectation by \ref{it:app:rec:rare}, and returns one
label with its re-isolated core representation, of the same length bound. The
proof record contains the oracle answers, retained lists, label values, and the
solver's computation; its expected size obeys the work bound. For quadratic
$F_0$ replace \cref{thm:int:solver} by \ref{it:app:rec:faces} on the core box
($3^k\poly(I)$, rational output) and take $B=\max\{2,N_Z3^k\}$.
\end{proof}

\begin{proof}[Proof of \cref{cor:int:native-implicit}]
Use the base choices above with two changes. Let
\[
 N'=\max\{1,\,k3^kN_Z\max(1,d-1)^k\},\qquad\tau=\frac{\rho\sigma}{2N'},
\]
and require additionally
\[
 2^{-J}\le\min\Bigl\{\frac\tau{8K_2A_0},\ \frac{g_0}{4K_3A_0}\Bigr\},
 \qquad M\ge\frac{2N'}\rho,
\]
where $K_2,K_3$ are as in \eqref{eq:rec:K} over the real box.
After \eqref{eq:int:label} passes at level $j$, apply the sign tests of
\cref{lem:rec:closure}(a) to the core coordinates of $f=F_\gamma(\cdot,z_j)$ on
$Q=\mathcal H_j$, which contains the core of every global minimizer by
\cref{lem:int:label}, and then the Hessian test of \cref{lem:rec:closure}(b) with
modulus $g_0$ on the remaining core box $Q'$. If both pass, return $z_j$ and the
patch $(Q',g_0)$; its minimizer is the global optimizer. Otherwise continue; the
fallback is that of \cref{thm:int:native}.

\emph{The extra event.} Fix a label $w$ in the bounding box, a face of the core
cube, and a core coordinate $i$ fixed on it, and condition on all coefficients
but $\gamma_i$. The free core stationary system of $F_\gamma(\cdot,w)$ on the face
does not involve $\gamma_i$ and has at most $\max(1,d-1)^k$ nonsingular solutions
(\ref{it:app:rec:bezout}); at each, the core derivative lies in $[-\tau,\tau]$
with probability at most $\tau/\sigma+1/M$. On $g_*>0$ the optimal core is a
nonsingular solution for $w=a_R$ on its face, by two-sided growth. Hence the
probability that $g_*\ge g_0$ and some active core derivative has magnitude at
most $\tau$ is at most $N'(\tau/\sigma+1/M)\le\rho$.

\emph{Stopping.} Outside both events, at level $J$ the label test passes
(\cref{lem:int:native-stop}), $\mathcal H_J\subseteq a+[-A_0h,A_0h]^k$, and the
midpoint test of a core coordinate at a bound has value at least
$\tau-2K_2A_0h\ge\frac34\tau$; free core coordinates have zero derivative at $a$
and are not fixed. Growth with the label fixed gives
$\nabla^2_{\mathcal U\mathcal U}f(a)\succeq2g_0I$ on the free core coordinates, and
the Hessian test value is at least $g_0-2K_3A_0h\ge g_0/2$. The failure probability
is at most $2\rho=1/(2B)$, there is no $c_d^k$ completion term, and a $q$-bit
evaluation uses \ref{it:app:rec:gls}.
\end{proof}

\begin{proof}[Proof of \cref{prop:int:hs}]
Write $Az=b$ as $Az\ge b$, $-Az\ge-b$; with the bound rows, the constraint matrix
remains totally unimodular, and tightening bounds keeps it so. At a rational core
point $v$ and rational linear coefficients, each $t\mapsto f_i(v,t)+\gamma_it$ is
a convex polynomial on $[\ell_i,u_i]$; extending it beyond each endpoint by the
tangent line there gives a convex function on $\R$ that agrees on the interval and
has rational values of polynomial length at rational points. The proximity-scaling
algorithm of \citet[Algorithm~4.2 and Theorem~4.3]{hochbaum1990-convex-separable-optimization-is-not} solves
separable convex integer minimization over totally unimodular systems exactly,
using a number of linear programs logarithmic in the numerical bounds, each with an
optimal extreme point computed exactly in polynomial bit time
\citep{grotschel1988-geometric-algorithms-and-combinatorial-optimization}, and function evaluations at prescribed
points. All data have polynomial length, so the oracle runs in polynomial bit
time, and the exact value is the evaluation at the returned integer point.
Infeasibility is detected by the linear program of the first stage.
\end{proof}

\begin{proof}[Proof of \cref{lem:int:potentials}]
\emph{Sufficiency.} Let $z'\in Y$ and $\Delta=z'-z$. It is a circulation; route
$\Delta_a>0$ units on the forward residual copy of $a$ and $-\Delta_a$ units on the
backward copy when $\Delta_a<0$; both copies are available because $z'$ satisfies
the bounds. By discrete convexity, $f_a(z'_a)-f_a(z_a)$ is at least the residual
arc cost times the number of units. Summing, $F(z')-F(z)$ is at least the cost of
this nonnegative residual circulation, which equals its reduced cost because
potential terms cancel around a circulation, and so is nonnegative.

\emph{Necessity.} A simple residual cycle using at most one copy of each
original arc admits one unit of feasible augmentation. Its objective change
is exactly the sum of its unit marginal costs, which is nonnegative by
optimality. This includes two-cycles on distinct arcs and one-arc loops.
The only simple cycle using both residual copies of an original arc is
their two-cycle, whose cost is nonnegative by convexity. Thus no simple
residual cycle has negative cost. Hence
shortest-path distances $\pi$ from an added source joined to every node by zero-cost
arcs are finite and satisfy $c_e+\pi_p-\pi_{p'}\ge0$ on every residual arc from $p$
to $p'$. A shortest path uses at most $s$ arcs. When unit marginal costs are
rational of bit length at most $H_m$, its length, and hence each potential,
has bit length polynomial in $s$ and $H_m$.
\end{proof}

\subsection{Flows under core-only noise}\label{app:int:flows}

\begin{proof}[Proof of \cref{cor:int:grid}]
The feasible residual set is fixed, so \cref{lem:rec:value}(b) and
\cref{lem:count:rounding} apply to the conditional value $V$. Every core
point belongs to a grid cell of side $h=1/m$, and a corner of that cell
has conditional value at most $V(v)+kLh^2/8$. Thus
$\min_{G_m}\ell-kLh^2/8\le\min V=f^*$. If $v_0$ minimizes the lower
answers on $G_m$, then the least upper answer satisfies
$F_\gamma(\hat v,\hat z)\le u(v_0)\le\ell(v_0)+\eta$; subtracting the
displayed lower bound proves the interval and gap claims. Feasibility
comes from the returned completion. The choice of $m$ makes $kLh^2/8$
at most $\varepsilon$. For integer flows the exact oracle of
\cref{prop:int:hs} has $\eta=0$; approximate calls have the separately
specified allowance $\eta$. No residual variable is rounded.
\end{proof}

We treat networks and totally unimodular systems together. Let
$Y=\{z\in\Z^{n_R}:Az=b,\ \ell\le z\le u\}$ with $A$ totally unimodular; a network
is the case of a node--arc incidence matrix. The objective is
$F_\gamma(v,z)=\phi(v)+\sum_if_i(v,z_i)+\gamma^{\mathsf T}v$ with $f_i(v,\cdot)$
convex on $[\ell_i,u_i]$. A \emph{multiplier map} is a polynomial vector
$\psi(w)\in\R^{n_R}$ such that $\psi(w)^{\mathsf T}z$ is constant on $Y$ for every
$w$; for a network $\psi_a(w)=\pi_{p'}(w)-\pi_p(w)$ for an arc $a=(p,p')$, and for a
totally unimodular system $\psi(w)=A^{\mathsf T}\lambda(w)$. The \emph{adjusted
costs} are $g_i(w,t)=f_i(w,t)-\psi_i(w)t$, and
\begin{equation}\label{eq:app:int:adjusted}
 F_\gamma(w,z)-F_\gamma(w,z')=\sum_i\bigl[g_i(w,z_i)-g_i(w,z'_i)\bigr]
 \qquad(z,z'\in Y).
\end{equation}

\begin{definition}[Charts]\label{def:int:charts}
For a network, a \emph{chart} is determined by a feasible flow $z$, a spanning
arborescence $\mathcal T$ rooted at an added source in the residual network of
$z$ augmented by zero-cost source arcs, an arc $a$, and an integer $t$ with
$t,t+1\in[\ell_a,u_a]$. Its potentials $\pi^{\mathcal T}(v)$ are the sums of the
residual unit marginals of $z$ along the tree paths, and its \emph{chart
polynomial} is $m(v)=f_a(v,t+1)-f_a(v,t)+\pi^{\mathcal T}_p(v)-\pi^{\mathcal T}_{p'}(v)$.
The family $\mathcal C$ of all charts has at most
$N_Z(2n_R+s)^sR_1$ members, $R_1=\sum_a(u_a-\ell_a)$.
\end{definition}

A unit difference $f_a(v,t+1)-f_a(v,t)$ has core degree at most $d-1$, so chart
polynomials and potentials have core degree at most $d-1$, adjusted costs have
total degree at most $d$, and all chart coefficients have bit length at most a
base bound $H_0=\poly_d(I)$: labels have base length, and tree paths have at most
$s$ arcs. A tree chosen by the algorithm is a member of $\mathcal C$, because the
family contains every structurally valid arborescence of every label. The
residual reduced costs of a flow under the potentials of such a tree are, up to
sign, chart polynomials.

\begin{proof}[Proof of \cref{thm:int:flow-interior}]
\emph{Algorithm.} Run the core search in exact mode with the convex flow oracle.
At level $j$, at the incumbent corner $c_j$ with flow $z_j$, compute shortest-path
distances in the residual network at $c_j$, choose a spanning arborescence of the
subgraph of tight arcs (zero reduced cost) by graph search, and form its potentials
$\pi(v)$. For every arc of the residual network of $z_j$ decide exactly whether
its reduced cost is nonnegative on $\mathcal H_j$, by minimizing it with
\cref{thm:int:solver} and comparing the algebraic minimum with zero; the source
arcs serve only to build the tree and are not tested. If all pass, \cref{lem:int:potentials}
shows that $z_j$ is optimal at every point of $\mathcal H_j$, which contains every
optimal core; so the slice $[0,1]^k\times\{z_j\}$ contains a global minimizer, and
\cref{thm:int:solver} solves it. After level $J$ use label enumeration as in
\cref{thm:int:native}. Every return is correct, with or without the interiority
premise.

\emph{Exceptional set.} Let $S\subseteq[0,1]^k$ be the union of the zero sets of
the chart polynomials that are not identically zero, and
$\mathcal E=\bigcup_{w}[-\nabla F_w](S)$ over feasible labels $w$, where
$F_w(v)=\phi(v)+\sum_af_a(v,w_a)$. Both are compact of dimension at most $k-1$. For
one chart polynomial $m$ and one label $w$, the set $\{\beta:\exists x\in[0,1]^k,\
m(x)=0,\ \beta+\nabla F_w(x)=0\}$ is defined by a formula with one block of $k$
variables, $k$ free variables, $O(k)$ atoms, and degree at most $\max\{d,2\}$. By
\ref{it:app:rec:qe} it is a Boolean combination of at most $\Xi_d(k)$ atoms of degree
at most $\Xi_d(k)$, with $\Xi_d(k)=(k+1)^{O_d(k^2)}$; as in the proof of
\cref{prop:rec:tube}(b), the product of its nonconstant atoms is a nonzero
polynomial of degree at most $\Xi_d(k)^3$ vanishing on it. So $\mathcal E$ lies in the
zero set of a nonzero polynomial of degree at most
$D_{\rm int}=\max\{1,N_Z\abs{\mathcal C}\Xi_d(k)^3\}$, and $\log D_{\rm int}\le
\poly_d(I)$. Nothing of this is computed by the algorithm.

\emph{Base choices.} Let $B$ be the label-enumeration budget of
\cref{thm:int:native}, $H_2\ge1$ a monomial bound on $\norm{\nabla^2F_w}$ over the
core box for all labels in the bounding box, $C_{\rm tube}=C(k,D_{\rm int})$ with
$C$ from \cref{prop:rec:tube}, and $C_{\rm growth}$ the constant of
\ref{it:app:rec:finite} for the reduced objective $V_0$ (core perturbed, flows as
the eliminated block). Choose $\rho=1/(4B)$, $g_0=\rho\sigma/(2k)$,
$A_0=2+kL/g_0$, $\delta=\rho\sigma/(4C_{\rm tube}H_2)$, the least $J$ with
$2^{-J}\le\delta/(2kA_0)$, and the least power of two
$M\ge\max\{2,2^J,4kC_{\rm growth}/\rho,4kC_{\rm tube}/\rho\}$.

\emph{Probability.} Let $a$ be an optimal core and $w$ an optimal flow at $a$. Then
$F_w+\gamma^{\mathsf T}v\ge V_\gamma$ with equality at $a$, so $a$ minimizes this
smooth function on the box; by the interiority premise $\gamma=-\nabla F_w(a)$, and
$\dist(\gamma,\mathcal E)\le H_2\dist(a,S)$. By \ref{it:app:rec:finite} and the
grid tube bound in the proof of \cref{prop:rec:tube}(c), the probability that the
projected growth is below $g_0$ or that $\dist(\gamma,\mathcal E)\le H_2\delta$ is at
most $\rho+C_{\rm tube}(H_2\delta/\sigma+\sqrt k/M)\le\frac32\rho<1/(2B)$.

\emph{Stopping.} Outside this event the optimal core $a$ is unique, as in
\cref{lem:int:native-stop} every point of $\mathcal H_J$ is within Euclidean distance
$kA_0h_J\le\delta/2$ of $a$, and $\dist(a,S)>\delta$; so $\mathcal H_J\cap S=\emptyset$.
Each reduced cost of the computed tree is, up to sign, a chart polynomial. If it is
identically zero it passes. Otherwise it has no zero on the connected box
$\mathcal H_J$, so it has constant sign there; it is nonnegative and, since
$c_J\notin S$, nonzero at $c_J$; hence it is positive on $\mathcal H_J$. The test
passes by level $J$.

\emph{Work.} Each level makes at most $2n_R$ exact box minimizations, of cost
$c_d^k\poly_d(I)$ in total over the $J$ levels after absorbing $J$ and $n_R$, plus
shortest paths and the core search. With one completion and the expected fallback,
the expected work is $[8^kQ_{\rm ex}+c_d^k]\poly_d(I)$. Quadratic objectives use
\ref{it:app:rec:faces} for the box minimizations and the completion.
\end{proof}

\begin{proof}[Proof of \cref{lem:int:proximity}]
We prove it for totally unimodular $A$; networks are a special case. Let $z\in Y$,
choose $\bar z\in Y_0$ nearest to $z$ in the $\ell_1$ norm, and put $x=z-\bar z\in
\ker A\cap\Z^{n_R}$. \emph{Conformal circuits.} If $x\ne0$, choose $y\in\ker A$
conformal to $x$ (same sign wherever $y_i\ne0$) with minimal nonempty support $S$. If
the vectors of $\ker A$ supported in $S$ formed a space of dimension at least two,
adding a small multiple of a nonparallel one and following until a coordinate
vanishes would give a conformal kernel vector with smaller support. So that space
is a line. Its primitive generator is given by signed maximal minors of the columns
in $S$, which lie in $\{0,\pm1\}$ by total unimodularity, so a conformal circuit
$c\in\{0,\pm1\}^{n_R}$ with support $S$ exists. Subtracting the largest integer
multiple $\min_{i\in S}\abs{x_i}$ of $c$ keeps $x$ integral and conformal and removes
a coordinate. Repeating gives $x=\sum_j\theta_jc^{(j)}$ with positive integers
$\theta_j$ and conformal circuits $c^{(j)}$.

\emph{Charging.} Each $\bar z+c^{(j)}$ satisfies $Az=b$ and lies between $\bar z$
and $z$ coordinatewise, hence within the bounds. If it also had all coordinates in
the intervals $I_i$, it would belong to $Y_0$ and be closer to $z$. So some
coordinate $i_j$ has $\bar z_{i_j}\in I_{i_j}$ but $\bar z_{i_j}+c^{(j)}_{i_j}\notin
I_{i_j}$: $\bar z_{i_j}$ is the endpoint of $I_{i_j}$ toward $z_{i_j}$, and
$\abs{x_{i_j}}=\dist(z_{i_j},I_{i_j})$. By conformality the multiplicities charged
to one coordinate $i$ sum to at most $\abs{x_i}$. Each circuit has at most $n_R$
nonzero entries, so $\norm x_1\le n_R\sum_j\theta_j\le n_R\sum_i\dist(z_i,I_i)$.
\end{proof}

\begin{lemma}[Optimal-face certificate, general form]\label{lem:int:face-general}
Let $\mathcal A$, $\mathcal D$, $\mathcal D_{\mathcal A}$, $c$, $\varsigma_i$, $T_1$,
$T_\infty$, $R_{\mathcal A}$, $H$, $K$ be as in \cref{thm:int:face}, with $K$ bounding
$\abs{\partial_{v_l}\partial_{z_i}f_i}$. Let $z^0\in Y$ be optimal at $c$ and let
$\psi$ be a polynomial multiplier map such that every $g_i(c,\cdot)$ is minimized over
the integers of $[\ell_i,u_i]$ at $z^0_i$. Let $I_i$ be the integer minimizers of
$g_i(c,\cdot)$, $Y_0=\{z\in Y:z_i\in I_i\}$, and assume (a)--(c) of
\cref{thm:int:face} for these adjusted costs. Then the conclusions of
\cref{thm:int:face} hold.
\end{lemma}

\begin{proof}
By \eqref{eq:app:int:adjusted}, $Y_0$ is the optimal set at $c$. For
$w\in\mathcal D_{\mathcal A}$, each $g_i(w,\cdot)$ is convex, constant on $I_i$ by (a),
and has first outside differences at least $\mu_{\rm out}$ by (b), so
$g_i(w,t)-g_i(w,z^0_i)\ge\mu_{\rm out}\dist(t,I_i)$ for integers $t$ in the native
interval. With \eqref{eq:app:int:adjusted},
\begin{equation}\label{eq:app:int:outside}
 F_\gamma(w,z)-V(w)\ge\mu_{\rm out}\,v_I(z),\qquad
 v_I(z)=\sum_i\dist(z_i,I_i),
\end{equation}
and every flow of $Y_0$ is optimal at $w$. Now let $v\in\mathcal D$ with face
projection $w$, write $v=w+\sum_i\varsigma_id_ie_i$ with $d\ge0$ supported on the normal
coordinates, and let $z\in Y$. Choose $\bar z\in Y_0$ by \cref{lem:int:proximity}. By the
derivative bounds,
$\varsigma_i\partial_{v_i}F_\gamma(w,z)\ge\varsigma_i\partial_{v_i}F_\gamma(c,\bar z)
-HR_{\mathcal A}-K\norm{z-\bar z}_1\ge\beta-HR_{\mathcal A}-Kn_Rv_I(z)$, and Taylor's
theorem in the core gives
\[
 F_\gamma(v,z)-V(w)\ge(\mu_{\rm out}-Kn_R\norm d_1)v_I(z)
 +(\beta-HR_{\mathcal A})\norm d_1-\frac H2\norm d_2^2
 \ge\Bigl(\beta-HR_{\mathcal A}-\frac H2T_\infty\Bigr)\norm d_1,
\]
using $\norm d_1\le T_1$, $\norm d_2^2\le T_\infty\norm d_1$, and (b). The coefficient
is positive by (c), so off the face $F_\gamma(v,z)>V(w)$, the value of a feasible
point on the face. If $\mathcal D$ contains every optimal core, every optimal core
lies in $\mathcal D_{\mathcal A}$, and $z^0\in Y_0$ is optimal there, so the slice of
$z^0$ contains a global minimizer.
\end{proof}

\begin{proof}[Proof of \cref{thm:int:face}]
For a network take $\psi_a(w)=\pi_{p'}(w)-\pi_p(w)$ from the tree potentials. At $c$
they certify optimality (\cref{lem:int:potentials}), which says exactly that each
$g_a(c,\cdot)$ is minimized at $z^0_a$, and $\psi(w)^{\mathsf T}z$ is constant on $Y$
because supplies are fixed. The intervals $I_a$ are found by binary search on the
nondecreasing differences $g_a(c,t+1)-g_a(c,t)$. For (c): on an interval with at
least three integers, convexity and equal values at three of its points make
$g_a(c,\cdot)$ constant on its convex hull, so $f_a(c,\cdot)$ is affine there and
$\partial_{tt}f_a(c,t)=0$; since $\partial_{tt}f_a(c+h\varsigma_ie_i,t)\ge0$ for small
$h\ge0$, the inward derivative $\varsigma_i\partial_{v_i}f_a(c,\cdot)$ is convex on that
interval. With linear interpolation on two-point intervals and substitution on
singletons, and tangent extensions, $\beta_i$ is the optimum of a separable convex
flow problem on the tightened bounds (\cref{prop:int:hs}). The polynomial identities
in (a) need $d+1$ integer representatives per interval because the adjusted cost has
total degree at most $d$. \Cref{lem:int:face-general} gives the conclusions.
\end{proof}

\begin{lemma}[Value margins away from zeros]\label{lem:int:margin}
Let $p\ne0$ be a rational polynomial of degree at most $d$ in $q\le k$ variables with
coefficients of bit length at most $H_0$, let $\delta>0$ be rational, let $Z_p$ be its
zero set in $[0,1]^q$, and let $\mathcal K=\{x\in[0,1]^q:\dist(x,Z_p)\ge\delta\}$ (the
whole cube if $Z_p=\emptyset$). If $\mathcal K\ne\emptyset$, then
$\min_{\mathcal K}\abs p\ge2^{-(B_0+1)}$ with $B_0\le\Xi'_d(k)\poly_d(H_0+\langle\delta\rangle+1)$
and $\Xi'_d(k)=(k+1)^{O_d(k^2)}$.
\end{lemma}

\begin{proof}
For $q=0$ the value is a nonzero rational with denominator at most $2^{H_0}$. Let
$q\ge1$. By compactness $m=\min_{\mathcal K}\abs p>0$. The formula
\[
 \exists x\ \forall y:\ x\in[0,1]^q\wedge\bigl[\neg(y\in[0,1]^q\wedge p(y)=0)\vee
 \norm{x-y}^2\ge\delta^2\bigr]\wedge0<p(x)^2<t
\]
has truth set $(m^2,\infty)$: a minimizer witnesses every $t>m^2$, and a witness lies
in $\mathcal K$, so no $t\le m^2$ is witnessed. It has two blocks of $q$ variables, one
free variable, $O(q)$ atoms, and degree at most $2\max\{d,1\}$. Clearing denominators
gives integer coefficients of bit length at most $\binom{q+d}d\poly(H_0+\langle\delta\rangle)$.
By \ref{it:app:rec:qe}, an equivalent quantifier-free formula has integer coefficients
of bit length at most $B_0$ of the stated form. The boundary point $m^2$ is a root of a
nonconstant atom polynomial, since otherwise all atom signs would be constant near
$m^2$. Remove a power of $t$ from it; the remaining integer polynomial has nonzero
constant term and coefficients of absolute value at most $2^{B_0}$, so Cauchy's bound
gives $m^2\ge1/(1+2^{B_0})\ge2^{-(B_0+1)}$, and $m\ge2^{-(B_0+1)}$.
\end{proof}

\begin{lemma}[Normal margins]\label{lem:int:normal}
For a face $\mathcal A$ of $[0,1]^k$, let $a_{\mathcal A}$ be the lexicographically
least minimizer of $V_\gamma$ on $\mathcal A$, $Y^*$ the set of flows optimal at
$a_{\mathcal A}$, and, for a normal coordinate $i$ of $\mathcal A$,
$\beta^*_{\mathcal A,i}=\min_{z\in Y^*}\varsigma_i\partial_{v_i}F_\gamma(a_{\mathcal A},z)$.
Under core-only noise with law $U_{\sigma,M}$,
$\Pr\{\exists\mathcal A,i:\beta^*_{\mathcal A,i}\in[0,\tau]\}\le k3^k(\tau/\sigma+1/M)$.
If the optimal core $a$ is unique and lies in the relative interior of $\mathcal A$,
then $a_{\mathcal A}=a$ and $\beta^*_{\mathcal A,i}\ge0$ for every normal $i$.
\end{lemma}

\begin{proof}
On $\mathcal A$, $V_\gamma$ equals $V_0$ plus the free noise terms plus a constant
given by the normal noises, so $a_{\mathcal A}$ depends only on the free coefficients;
$Y^*$ depends only on $a_{\mathcal A}$, because core noise is constant in $z$. Hence
$\beta^*_{\mathcal A,i}=\bar b+\varsigma_i\gamma_i$ with $\bar b$ independent of
$\gamma_i$ given the other coefficients, and \eqref{eq:model:interval} bounds the conditional
probability by $\tau/(2\sigma)+1/M$. A union over at most $k3^k$ pairs gives the
bound. If $a$ is the unique optimal core and lies on $\mathcal A$, every minimizer on
$\mathcal A$ is a global one, so $a_{\mathcal A}=a$. For $z\in Y^*$,
$F_\gamma(\cdot,z)\ge V_\gamma\ge V_\gamma(a)=F_\gamma(a,z)$, so $a$ minimizes the smooth
function $F_\gamma(\cdot,z)$ on the box, and its inward derivatives are nonnegative.
\end{proof}

\begin{proof}[Proof of \cref{thm:int:flow-boundary}]
\emph{Algorithm.} At each level $j\le J$ and for each face $\mathcal A$ whose fixed
bounds lie in the corresponding intervals of $\mathcal H_j$, let $\mathcal D=\mathcal H_j$,
$c$ the midpoint of $\mathcal D_{\mathcal A}$, compute an optimal flow, tree potentials,
and intervals at $c$, and test (a)--(c) of \cref{thm:int:face}; on success solve the
whole core slice by \cref{thm:int:solver}. After level $J$ enumerate labels. Correctness
follows from \cref{thm:int:face} and \cref{prop:rec:search}(a).

\emph{Analysis sets.} For a face $\mathcal A$ with $q\ge1$ free coordinates, let
$S_{\mathcal A}$ be the union of the zero sets in the face of the chart polynomials whose
restriction to the face is not identically zero, and
$\mathcal E_{\mathcal A}=\bigcup_w[-\nabla_{\mathcal A}F_w](S_{\mathcal A})$. As in the
interior case, $\mathcal E_{\mathcal A}$ lies in the zero set of a nonzero polynomial of
degree at most $D=\max\{1,N_Z\abs{\mathcal C}\Xi_d(k)^3\}$. Let $H\ge1$ bound core
Hessians and $K$ the mixed derivatives, $C_{\rm tot}=3^kC(k,D)$,
$N_{\rm nor}=k3^k$, and let $C_{\rm growth}$ and $B$ be as before. Choose
$\rho=1/(8B)$, $g_0=\rho\sigma/(2k)$, $A_0=2+kL/g_0$,
$\delta=\min\{\frac14,\rho\sigma/(4C_{\rm tot}H)\}$, $\tau=\rho\sigma/(4N_{\rm nor})$, and
$\mu_0=2^{-(B_0+1)}$ from \cref{lem:int:margin} with $\delta/2$, uniformly for all chart
polynomials restricted to all faces. Let $J$ be least with
\[
 2^{-J}\le\min\Bigl\{\frac\delta{4kA_0},\frac\tau{16kHA_0},
 \frac{\mu_0}{4k\max\{1,n_RK\}A_0}\Bigr\},
\]
and $M$ the least power of two with $M\ge\max\{2,2^J,4kC_{\rm growth}/\rho,
4kC_{\rm tot}/\rho,4N_{\rm nor}/\rho\}$. Then $J,\log_2M\le f_d(k)\poly_d(I)$, because
$\log(1/\mu_0)$ is of that form; no sampled quantity enters these choices.

\emph{Probability.} The growth event costs at most $\rho$. At an optimal core $a$ in
the relative interior of a face $\mathcal A$ with $q\ge1$, every optimal flow $w$ makes
$F_w|_{\mathcal A}+\gamma^{\mathsf T}v$ minimal at $a$ in the free directions, so
$\gamma_{\mathcal A}=-\nabla_{\mathcal A}F_w(a)$ and
$\dist(\gamma_{\mathcal A},\mathcal E_{\mathcal A})\le H\dist(a,S_{\mathcal A})$; the tube
event that this is at most $H\delta$ for some face costs at most
$C_{\rm tot}(H\delta/\sigma+k/M)\le\rho/2$. \Cref{lem:int:normal} costs at most
$N_{\rm nor}(\tau/\sigma+1/M)\le\rho/2$. The total is at most $2\rho=1/(4B)$.

\emph{Stopping.} Outside these events let $a$ be the unique optimal core and
$\mathcal A$ its face. At level $J$, $\mathcal H_J\subseteq a+[-A_0h,A_0h]^k$, so the face
is tried; its normal widths are at most $A_0h$, $T_1\le kA_0h$, $T_\infty\le A_0h$,
$R_{\mathcal A}\le2kA_0h$, and $\mathcal D_{\mathcal A}$, including $c$, lies within
Euclidean distance $kA_0h\le\delta/4$ of $a$. If $q\ge1$, then $\dist(a,S_{\mathcal A})>\delta$,
so $\mathcal D_{\mathcal A}$ is at distance more than $\delta/2$ from $S_{\mathcal A}$. A
within-interval unit difference of the computed adjusted costs is a chart polynomial
vanishing at $c\notin S_{\mathcal A}$, hence identically zero on the face: (a) passes.
A first outside difference is positive at $c$ because the interval is the full set of
minimizers, so it has no zero on the connected set $\mathcal D_{\mathcal A}$ and, by
\cref{lem:int:margin}, is at least $\mu_0>\mu_0/4\ge n_RKT_1$ there: (b) passes. For a
vertex face the same holds with constant charts. Consequently $Y_0$ is the optimal set
at every point of $\mathcal D_{\mathcal A}$, in particular at $a$, so
$\beta_i\ge\beta^*_{\mathcal A,i}-HkA_0h>\tau-HkA_0h$ by \cref{lem:int:normal}, and
$\beta-HR_{\mathcal A}-\frac H2T_\infty>\tau-4kHA_0h\ge\frac34\tau$: (c) passes.

\emph{Work and output.} At most $3^k$ face attempts per level, each with polynomially
many flow calls, interval searches, identities, and $c_d^k$-cost box minimizations on
data of length $f_d(k)\poly_d(I)$, give with the core search and the expected fallback
the bound $f_d(k)Q_{\rm ex}\poly_d(I)$. The completion and the fallback return one flow
and one common-root core representation of length $f_d(k)\poly_d(I)$.
\end{proof}

\begin{lemma}[Affine margins]\label{lem:int:affine-margin}
Let $p(x)=p_0+\pi^{\mathsf T}x$ be a nonzero affine polynomial on $[0,1]^q$ whose
coefficients have numerators and denominators of bit length at most $H_0$, and let
$Z_p$ be its zero set in the cube. If $Z_p\ne\emptyset$, then
$\abs{p(x)}\ge2^{-H_0}\dist(x,Z_p)$; otherwise $\abs p\ge2^{-(q+1)H_0}$ on the cube.
\end{lemma}

\begin{proof}
Let $\pi_{\min}\ge2^{-H_0}$ be the least nonzero $\abs{\pi_i}$. If $Z_p\ne\emptyset$ and
$p(x)>0$, some vertex has $p\le0$. Move from $x$ toward it one coordinate with
$\pi_i\ne0$ at a time; $p$ decreases at rate at least $\pi_{\min}$ per unit length and
reaches zero after an $\ell_1$ path of length at most $p(x)/\pi_{\min}$, which bounds the
Euclidean distance to $Z_p$. Negative values are symmetric. If $Z_p=\emptyset$, $p$
has one sign on the cube and its least absolute value is attained at a vertex, where
it is a nonzero rational whose denominator divides a product of at most $q+1$
denominators.
\end{proof}

\begin{proof}[Proof of \cref{cor:int:bilinear}]
All unit marginals $\psi_a(t+1)-\psi_a(t)+(W^{\mathsf T}v)_a$, all potentials, and all
chart polynomials are affine in the core with base-bounded coefficients. Replace
$\mu_0$ in the proof of \cref{thm:int:flow-boundary} by
$2^{-(k+1)H_0}\delta/2$, valid by \cref{lem:int:affine-margin} at distance $\delta/2$;
it has polynomial length, so $J,\log_2M\le\poly_d(I)$. The identities in (a) are
coefficient comparisons, the minima in (b) are endpoint choices, and the inward
derivative $\varsigma_i(\partial_i\phi(c)+\gamma_i+(Wz)_i)$ is linear in $z$, so $\beta_i$
is a linear-cost flow problem on the tightened bounds. The tube degree bound still uses
\ref{it:app:rec:qe} for the possibly nonlinear $\phi$, but only its logarithm enters. The
per-level work is polynomial and is absorbed by $8^kQ_{\rm ex}$; the completion
minimizes $\phi$ plus a linear function, at cost $c_d^k\poly_d(I)$; quadratic $\phi$
gives rational output by \ref{it:app:rec:faces}.
\end{proof}

\subsection{Totally unimodular systems}\label{app:int:tu}

Substitute fixed coordinates, check consistency, and delete dependent rows; the
remaining $A$ is totally unimodular with full row rank $m$, and every native interval
has positive length.

\begin{lemma}[A compact dual certificate]\label{lem:int:tu-dual}
Let $z^0$ be an integer optimum at a rational core point $c$. There is
$\lambda\in\R^m$ with
\begin{equation}\label{eq:app:int:slopes}
 f_i(c,z^0_i)-f_i(c,z^0_i-1)\le(A^{\mathsf T}\lambda)_i\ \ (z^0_i>\ell_i),\qquad
 (A^{\mathsf T}\lambda)_i\le f_i(c,z^0_i+1)-f_i(c,z^0_i)\ \ (z^0_i<u_i).
\end{equation}
Every solution certifies optimality of $z^0$. With a rational bound $\bar V\ge1$ on the
absolute values of all unit marginals over the core cube, a solution exists in
$[-R,R]^m$, $R=m\bar V+1$, and a vertex of \eqref{eq:app:int:slopes} intersected with this
box is computed in polynomial time; its active basis $\mathsf C$, a nonsingular
submatrix of $A^{\mathsf T}$ with unit rows appended, has an inverse with entries in
$\{0,\pm1\}$.
\end{lemma}

\begin{proof}
Let $\hat f_i$ be the piecewise-linear interpolation of $f_i(c,\cdot)$ at the integers of
$[\ell_i,u_i]$; it is convex. On each integer unit cell the function
$\sum_i\hat f_i$ is affine and the cell intersected with $\{Az=b\}$ is a polytope with a
totally unimodular matrix and integral right-hand side, so its vertices are integral.
Hence the minimum of $\sum_i\hat f_i$ over the real polytope equals the integer minimum,
and $z^0$ is a real optimum. Linear programming duality for the epigraph formulation
gives $\lambda$ such that $z^0$ minimizes $\sum_i[\hat f_i(z_i)-(A^{\mathsf T}\lambda)_iz_i]$
over the box; for convex piecewise-linear functions with integer breakpoints this is
\eqref{eq:app:int:slopes}. Conversely, \eqref{eq:app:int:slopes} and discrete convexity
make each $f_i(c,t)-(A^{\mathsf T}\lambda)_it$ minimal at $z^0_i$, and
$\lambda^{\mathsf T}Az=\lambda^{\mathsf T}b$ on $Y$.

Every column has an inequality, and a line in the solution set would be orthogonal to
all columns, which full row rank excludes; so the nonempty polyhedron has a vertex,
determined by $m$ independent active rows forming a nonsingular signed submatrix of
$A^{\mathsf T}$, whose inverse has entries in $\{0,\pm1\}$ by Cramer's rule. Its
coordinates are at most $m\bar V$ in absolute value, so the box contains it. A vertex of
the bounded polyhedron is found by exact linear programming, and its active rows are
extracted from the original inequalities and box rows; appending unit rows to
$A^{\mathsf T}$ preserves total unimodularity.
\end{proof}

\begin{proof}[Proof of \cref{thm:int:tu} and \cref{cor:int:tu-ineq}]
At a core point the algorithm uses the totally unimodular oracle of \cref{prop:int:hs},
computes $\lambda$ and its basis $\mathsf C$ by \cref{lem:int:tu-dual}, and keeps the basis
symbolically: $\lambda(v)=\mathsf C^{-1}h(v)$, where $h(v)$ consists of unit marginals and
the constants $\pm R$. This is a multiplier map of core degree at most $d-1$ with
base-bounded coefficients, and the adjusted costs at $c$ are minimized at $z^0$. The chart
family consists of the polynomials $f_i(v,t+1)-f_i(v,t)-(A^{\mathsf T}\lambda(v))_i$ over all
labels, ordered bases of at most $2n_R+2m$ candidate rows, coordinates, and adjacent
pairs; it has at most $N_Z(2n_R+2m)^mR_1$ members, so all logarithms remain polynomial.
Only the original adjusted marginals are tested on the hull; the artificial box rows
serve to select a bounded vertex at $c$. \Cref{lem:int:proximity,lem:int:face-general}
hold for $A$, and the proofs of \cref{thm:int:flow-boundary} and \cref{cor:int:bilinear}
apply verbatim with these charts. For \cref{cor:int:tu-ineq}, the slack system
$Cz+s=b$, $0\le s_i\le b_i-\min_{\ell\le z\le u}C_iz$, has the totally unimodular matrix
$[C\ I]$ and projects bijectively onto $Y$; slacks have zero cost and coupling, so the
objective, its core derivatives, and $L$ are unchanged, and the added bounds have
polynomial length.
\end{proof}

\subsection{Strong fields}\label{app:int:strong-field}

\begin{proof}[Proof of \cref{thm:int:strong-field}]
\emph{Enclosures.} For a monomial derivative, the exact range of each univariate power
on its interval is computed from the endpoints, adding $0$ for an even positive power
on an interval containing $0$; the product of these intervals, scaled by the
coefficient, and the sum over monomials enclose $\partial_iF_0$ on the hull. All
endpoints have polynomial length, and each $q_i$ is computed by counting the grid
indices of the closed interval $[-\overline\partial_i,-\underline\partial_i]$.

\emph{Pinning.} If $\gamma_i>-\underline\partial_i$, then $\partial_iF_\gamma>0$ on the
whole continuous hull; for any feasible $x$ with $x_i>\ell_i$, decreasing $x_i$ to
$\ell_i$, continuously or between integer labels, strictly decreases $F_\gamma$, and the
product domain keeps the point feasible. So every global minimizer has $x_i=\ell_i$; the
case $\gamma_i<-\overline\partial_i$ is symmetric. The bad event is
$\gamma_i\in[-\overline\partial_i,-\underline\partial_i]$, which leaves
coordinate $i$ unpinned by this test. These events depend on one coefficient
each and are independent with probabilities $q_i$.

\emph{Separation.} After substituting the pinned values, each remaining monomial has
pairwise adjacent bad variables, which lie in one component of the induced subgraph.
So $F_\gamma=c_\gamma+\sum_Cf_{C,\gamma}$ on a product of component domains, and choosing
a minimizer of each summand minimizes the sum. Each component is solved by
\cref{cor:int:mixed-solver}, or for quadratics by \ref{it:app:rec:faces}, in
$A(C)\poly_d(I+\log_2M)$ with $A(C)=\prod_{i\in C}a_i$; the separated data have length
$\poly_d(I+\log_2M)$.

\emph{Expectation.} A connected set of size $t$ containing a vertex $v$ has a rooted
ordered spanning tree with root $v$; there are at most $4^{t-1}$ shapes and at most
$\Delta_+^{t-1}$ choices of neighbors, so at most $(4\Delta_+)^{t-1}$ such sets. A
connected set $S$ is a bad component only if all its vertices are bad, with probability
at most $\prod_{i\in S}q_i$. Hence
\[
 \E\sum_{C}A(C)\le\sum_{S\text{ connected}}\prod_{i\in S}a_iq_i
 \le\sum_va_vq_v\sum_{t\ge1}(4\Delta_+\beta)^{t-1}=\frac{\sum_va_vq_v}{1-4\Delta_+\beta},
\]
which gives \eqref{eq:int:sf-work}. Under \eqref{eq:int:sf-regime},
\eqref{eq:model:interval} gives $q_i\le R_i/(2\sigma)+1/M$, so
$a_iq_i\le\frac1{16\Delta_+}+\frac1{16\Delta_+}$.

\emph{Refinement.} Put
$q'=q+\lceil\log_2(n+1)\rceil$. For each component separately, use
\cref{thm:int:solver} to obtain point error at most $2^{-q'}$ relative to
its stored optimizer, value width at most $2^{-q'}$, and a feasible
rational point with certified gap at most $2^{-q'}$. Clip continuous
approximations to their boxes and keep integer labels exact. Combining the
point approximations with pinned coordinates gives Euclidean error at most
$\sqrt n\,2^{-q'}\le2^{-q}$; adding value intervals and certified gaps
gives width and gap at most $n2^{-q'}\le2^{-q}$. The gap witnesses and the
point approximations may be refined separately; no distance conclusion is
inferred from an objective gap. This adds only $O(\log(n+1))$ bits. If all
coordinates are pinned, the rational point and value are returned directly.
The proof
record (enclosures, pinning comparisons, partition, component computations) is
checked within the same bound; independence and the noise condition enter only the
expectation.
\end{proof}

\subsection{The low-rank lattice theorem}\label{app:int:lowrank}

\begin{proof}[Proof of \cref{thm:int:lowrank}]
\emph{Expansion and scalar recourse.} Work with the expanded coefficients
of all listed pieces. Expanding a shifted or factored piece of degree at
most the fixed $d_*$ takes polynomial bit work and preserves polynomial
coefficient length. Let $Q_{\rm den}$ be the denominator product specified
in \cref{sec:int:lowrank}; all expanded coefficients have denominator
dividing it. For a rational $a\in\R^k$ put
$W(a)=\frac\alpha2\norm a^2+\min_{x\in X}[G(x)-\alpha a^{\mathsf T}Tx]$.
The inner problem separates into
$\psi_j(t)=g_j(t)-\alpha(T^{\mathsf T}a)_jt$ on
$\{L_j,\ldots,U_j\}$. Convexity gives
$2\psi_j(t+1)\le\psi_j(t)+\psi_j(t+2)$, so its forward differences are
nondecreasing even when consecutive arguments lie in different pieces.
For each difference, locate the piece at each argument by exact comparisons
with the rational knots and evaluate both polynomials. At a knot either
adjacent piece gives the same value by continuity.

Binary search for the first $t$ with
$\psi_j(t+1)-\psi_j(t)\ge0$ returns a minimizer; if none exists choose
$U_j$. A singleton interval is evaluated directly. The neighboring signs
at the returned point, when the neighbors exist, certify optimality by
summing the nondecreasing differences to any other integer. Ties,
constant pieces, affine pieces, and changes of derivative at knots are
allowed. There are $O(1+\log(U_j-L_j+1))$ comparisons per coordinate,
without enumerating its labels. This returns $W(a)$ and an attaining
integer witness $x(a)$. If $k=0$, these scalar searches solve the original
separable objective directly; henceforth assume $k\ge1$.

\emph{Square completion.} Since $\frac\alpha2\norm a^2-\alpha a^{\mathsf T}Tx
=\frac\alpha2\norm{a-Tx}^2-\frac\alpha2\norm{Tx}^2$, $W(a)=\min_x[F(x)+\frac\alpha2
\norm{a-Tx}^2]$, and $W-\frac\alpha2\norm a^2$ is a minimum of affine functions, so $W$
has upper coordinate curvature $\alpha$. With $Z_\xi(a)=W(a)+\xi^{\mathsf T}a$ and
$C_\xi=\norm\xi^2/(2\alpha)$,
\[
 Z_\xi(a)=\min_{x\in X}\Bigl[F_\xi(x)+\frac\alpha2\Bigl\lVert a-Tx+\frac\xi\alpha\Bigr\rVert^2\Bigr]-C_\xi .
\]
For an optimal $x^*$, $a=Tx^*-\xi/\alpha$ lies in the box
$\mathcal A=\prod_i[\ell_i-\sigma_i/\alpha,u_i+\sigma_i/\alpha]$, so
$\min_{\mathcal A}Z_\xi=F_\xi^*-C_\xi$, and every witness satisfies
$F_\xi(x(a))\le Z_\xi(a)+C_\xi$.

\emph{Search and count.} Run the corrected-corner search of \cref{thm:count:cells} on
$\mathcal A$ with $f=W$, $c=\xi$, curvature $\alpha$ in every coordinate, and the
isotropic balanced meshes of \cref{def:count:mesh} ($h_j=s2^{-j}$), through the
predetermined level $J=\log_2M$, with exact evaluations of $Z_\xi$. Since $m_{iJ}\le2^J=M$,
\cref{cor:count:levels} with the coordinatewise laws $U_{\sigma_i,M}$ gives expected
level counts at most $H_{\rm rat}$, and the expected number of evaluations is at most
$2^k+8^kH_{\rm rat}J$. The comparison in
\cref{cor:count:levels} applies only to coordinates with interior grid
nodes; coordinates not yet subdivided contribute their two endpoints.
Zero, dependent, or constant rows of $T$ cause no exception: each auxiliary
width is positive because $\sigma_i>0$.

\emph{Exact termination.} By \cref{thm:count:cells}(iii), the incumbent value satisfies
$U_J\le\min_{\mathcal A}Z_\xi+E_J$, and its witness has
$0\le F_\xi(x_{\rm inc})-F_\xi^*\le U_J+C_\xi-(\min_{\mathcal A}Z_\xi+C_\xi)\le E_J$, and
\begin{equation*}
 E_J=\frac\alpha8\sum_ih_{iJ}^2
 \le\frac{k\alpha s^2}{8M^2}
 \le\frac1{8D_0M}<\frac1{D_0(M-1)}.
\end{equation*}
Here $D_0=2Q_{\rm den}^3$, and the choice
$M\ge\max\{2,k\alpha s^2D_0\}$ gives the second inequality. Every base denominator divides
$Q_{\rm den}$. For integer $x$, $g_j(x_j)\in\frac1{Q_{\rm den}}\Z$,
$\frac\alpha2\norm{Tx}^2\in\frac1{2Q_{\rm den}^3}\Z$, and, writing
$\xi_i=\sigma_i(2K_i-(M-1))/(M-1)$ with the sampled index $K_i$,
$\xi^{\mathsf T}Tx\in\frac1{Q_{\rm den}^2(M-1)}\Z$. So every value $F_\xi(x)$ lies in
$\frac1{D_0(M-1)}\Z$. The optimum lies in this lattice because $X$ is finite.
The nonnegative incumbent gap is strictly smaller than its spacing, so it
is zero and $x_{\rm inc}$ is optimal; its value is
computed exactly. One common grid size for all rows is what makes the original values
share the single denominator factor $M-1$; the auxiliary values $Z_\xi$ and $C_\xi$ can
involve $(M-1)^2$, and no lattice property of them is used. An integer unary
monomial is an integer, whatever its degree, and a rational knot only
selects a piece; neither contributes another sampling denominator.

\emph{Bit lengths and work.} The sum of the expanded denominator lengths
bounds $\log Q_{\rm den}$, so $\log D_0=O(I)$. Row extrema are obtained by
rational endpoint sums. The widths, $s$, and $k\alpha s^2D_0$ therefore have
polynomial length. Comparing rationals computes the least adequate power
of two in polynomial work; it does not enumerate its support. Thus
$J=\log_2M\le\poly_{d_*}(I)$, sampling uses exactly $kJ$ fair bits, and
each $\xi_i$ has polynomial length.

Every feasible integer has input-bounded bit length. Raising a $B$-bit
integer to a power at most $d_*$ uses $O(d_*B)$ bits; rational Horner
evaluation on a listed piece has polynomial numerator and denominator
length. Piece lookup compares rational knots with binary integers by
integer products, and both lookup and the number of scalar comparisons
are polynomial in $I$. This also covers negative labels and knots within
a unit integer step. Continuity and convexity can be checked with rational
endpoint evaluations and univariate sign determination on each piece;
the input length counts the complete list.

At level $j\le J$, each queried auxiliary coordinate is its base endpoint
plus an integer index times $w_i/m_{ij}$, with polynomial length in $I+j$.
Consequently $\alpha(T^{\mathsf T}a)_j$, all scalar differences, $W(a)$,
$Z_\xi(a)$, $E_j$, and each original witness value have uniformly
polynomial length. Each query combines only polynomially many rationals;
denominators do not accumulate from unrelated calls because an incumbent
is selected by comparison and each new query is constructed from the base
data and its mesh indices. Cell addresses have $O(kJ)$ bits. Every query
and every cell therefore costs $\poly_{d_*}(I)$ with an exponent
independent of $k$, even on a draw visiting many cells. Multiplying this
cost by $2^k+8^kH_{\rm rat}J$ proves the expected bound. The final integer
vector has input-bounded length, and the same fixed-degree evaluation
bound gives a polynomial-length rational optimal value on every draw.
\end{proof}
